\documentclass[12pt]{article}
\usepackage{amsmath,amssymb,amsfonts,amsthm}
\usepackage[margin=1in]{geometry}

\begin{document}

\begin{center}
{\Large\bfseries Chapter 4, Problem 25}\\[6pt]
Byron \& Fuller, \textit{Mathematics of Classical and Quantum Physics}
\end{center}

\bigskip

\textbf{Problem.} Consider a four-dimensional real vector space with basis vectors $\{e_1, e_2, e_3, e_4\}$ and define a multiplication of basis vectors as follows:
\[
e_\mu e_\nu = e_\nu \quad (\mu = 1, 2, 3), \qquad e_4 e_\nu = e_\nu \quad (\nu = 1, 2, 3),
\]
\[
e_\mu e_4 = e_\mu \quad (\mu = 1, 2, 3), \qquad e_4 e_4 = e_4.
\]

Because of this last equation it is customary to write $e_4 = 1$ (identity), and hence to write a general vector in this space as
\[
q = e_1 q_1 + e_2 q_2 + e_3 q_3 + q_4,
\]
where $q_k$ are real scalars.

\begin{itemize}
\item[(a)] Extend the definition of vector multiplication to all vectors in the space by linearity. (Following Sir W.\ R.\ Hamilton, we call the elements of the space \textit{quaternions}.) Let $q = \alpha + \beta$ and $p = \gamma + \delta$ be any two quaternions, where $\alpha$ and $\gamma$ are ``pure quaternions'' (zero scalar part) and $\beta, \delta$ are scalars.
Show that $pq = -(\alpha,\gamma) + \beta\gamma + \delta\alpha + \alpha \times \gamma + \beta\delta$.
\item[(b)] When does $pq = 0$? That is, either $p = 0$ or $q = 0$.
\item[(c)] Show that the unit quaternion equation is equal to $1$; that is, $e_4 q = q$.
\item[(d)] Prove that $|pq| = |p||q|$. Hence the set of quaternions form what is called a \textit{normed} algebra.
\item[(e)] Show that the set of quaternions with unit norm (quaternions $q$ for which $|q| = 1$) form a group under quaternion multiplication.
\item[(f)] Let $q = \alpha + \beta$, $q^* = -\alpha + \beta$ (a quaternion conjugate); define its quaternion conjugate to be $q^* = -\alpha + \beta$. Show that $|q|^2 = qq^*$. This shows that the quaternions form a non-commutative field.
\item[(g)] Show that a realization of the quaternion field is obtained by writing
\[
j = \sqrt{-1},
\]
where $j$ acts on complex 2-vectors by...
\end{itemize}

\bigskip
\textbf{Solution.}

The quaternion basis elements satisfy: $e_1^2 = e_2^2 = e_3^2 = -1$ (from the multiplication table, where $e_\mu e_\mu = -e_4 = -1$ for $\mu = 1,2,3$), and the cross-product relations $e_1 e_2 = e_3$, $e_2 e_3 = e_1$, $e_3 e_1 = e_2$, with $e_\mu e_\nu = -e_\nu e_\mu$ for $\mu \neq \nu$.

We identify $e_1 = \mathbf{i}$, $e_2 = \mathbf{j}$, $e_3 = \mathbf{k}$, $e_4 = 1$.

\textbf{(a)} Let $p = \gamma + \delta$ and $q = \alpha + \beta$, where $\alpha = \sum_{i=1}^3 a_i e_i$, $\gamma = \sum_{i=1}^3 c_i e_i$ are pure quaternions and $\beta, \delta$ are real scalars.

\begin{align*}
pq &= (\gamma + \delta)(\alpha + \beta) = \gamma\alpha + \gamma\beta + \delta\alpha + \delta\beta.
\end{align*}

Now $\gamma\beta = \beta\gamma$ (scalar times pure quaternion commute), $\delta\alpha = \delta\alpha$, and $\delta\beta = \beta\delta$ (product of scalars).

For the product of two pure quaternions $\gamma\alpha$:
\[
\gamma\alpha = \sum_{i,j} c_i a_j\, e_i e_j.
\]
Using $e_i e_j = -\delta_{ij} + \epsilon_{ijk}e_k$:
\[
\gamma\alpha = -\sum_i c_i a_i + \sum_{i,j,k}\epsilon_{ijk}c_i a_j e_k = -(\gamma, \alpha) + \gamma \times \alpha.
\]

where $(\gamma,\alpha) = \sum_i c_i a_i$ is the Euclidean dot product and $\gamma \times \alpha$ is the cross product.

Therefore:
\[
\boxed{pq = -(\gamma,\alpha) + \beta\gamma + \delta\alpha + \gamma\times\alpha + \beta\delta}.
\]

\textbf{(b)} We show $pq = 0 \Rightarrow p = 0$ or $q = 0$ (no zero divisors).

Define $|q|^2 = q_1^2 + q_2^2 + q_3^2 + q_4^2$ for $q = \sum q_i e_i$. One can verify that $|pq|^2 = |p|^2|q|^2$ (see part (d)). If $pq = 0$, then $|p|^2|q|^2 = 0$, so $|p| = 0$ or $|q| = 0$, hence $p = 0$ or $q = 0$.

\textbf{(c)} $e_4 q = e_4(\sum_i q_i e_i) = \sum_i q_i(e_4 e_i) = \sum_i q_i e_i = q$. So $e_4 = 1$ acts as the identity. \checkmark

\textbf{(d)} For $q = \alpha + \beta$, define $\bar{q} = -\alpha + \beta$. Then:
\[
q\bar{q} = (\alpha+\beta)(-\alpha+\beta) = -\alpha^2 + \alpha\beta - \beta\alpha + \beta^2 = |\alpha|^2 + \beta^2 = |q|^2,
\]
since $\alpha^2 = -|\alpha|^2$ for a pure quaternion and $\alpha\beta = \beta\alpha$.

Similarly $|pq|^2 = (pq)\overline{(pq)} = (pq)(\bar{q}\bar{p}) = p(q\bar{q})\bar{p} = p|q|^2\bar{p} = |q|^2 p\bar{p} = |q|^2|p|^2$.

Therefore $|pq| = |p||q|$. \qed

\textbf{(e)} The set $S^3 = \{q : |q| = 1\}$ forms a group:
\begin{itemize}
\item \textbf{Closure:} $|pq| = |p||q| = 1\cdot 1 = 1$. \checkmark
\item \textbf{Identity:} $|1| = 1$. \checkmark
\item \textbf{Inverse:} $q^{-1} = \bar{q}/|q|^2 = \bar{q}$ when $|q| = 1$, and $|\bar{q}| = |q| = 1$. \checkmark
\item \textbf{Associativity:} Quaternion multiplication is associative. \checkmark
\end{itemize}

\textbf{(f)} As shown in part (d):
\[
q\bar{q} = |q|^2 = q_1^2 + q_2^2 + q_3^2 + q_4^2.
\]
Every nonzero quaternion has a multiplicative inverse $q^{-1} = \bar{q}/|q|^2$, and there are no zero divisors (part (b)). The quaternions thus form a division algebra (non-commutative field). \qed

\end{document}
